\documentclass[11pt]{article}
\usepackage[a4paper,margin=2.4cm]{geometry}
\usepackage{fontspec}
\usepackage{xeCJK}
\setmainfont{Noto Serif}
\setCJKmainfont{Microsoft JhengHei}
\setCJKsansfont{Microsoft JhengHei}
\setlength{\parindent}{2em}
\setlength{\parskip}{0.5em}
\begin{document}

\section*{36. 理想微分與不同}

\begin{center}
\emph{J. Ber. d. DMV 39 (1930), S. 17}
\end{center}

2. E. Noether，哥廷根：理想微分與不同。

本文說明，代數數域的不同如何可被理解為一個定義\emph{理想}的\emph{微分商}。這種定義與通常定義相一致，是由一些本身也很有意義的結構定理推出的；這些定理在類似意義下表示 Lagrange 插值公式的推廣。

詳細論述將發表於 \emph{Mathematische Annalen}。

\end{document}
